\begin{table}[htbp]\centering\small
\caption{Circunstancias locales del eclipse en \siteLat\textdegree{}\,N, \siteLon\textdegree{}\,E, \siteElevDEM\,m, calculadas de JPL DE440s con $\Delta T = \deltaT$\,s y $k = 0{,}2725076$. La altura solar incluye refracción.}
\label{tab:circ}
\begin{tabular}{lccS[table-format=-2.2]S[table-format=3.1]S[table-format=1.4]S[table-format=1.4]}
\toprule
Evento & UTC & CEST & {Altura (\textdegree)} & {Acimut (\textdegree)} & {Magnitud} & {Obscuración} \\
\midrule
Primer contacto (C1) & 17:35:27.6 & 19:35:27 & 14.70 & 277.0 & 0.0000 & 0.0000 \\
Segundo contacto (C2) & 18:29:25.6 & 20:29:25 & 4.85 & 285.6 & 1.0324 & 1.0000 \\
Máximo del eclipse & 18:29:59.9 & 20:29:59 & 4.75 & 285.7 & 1.0324 & 1.0000 \\
Tercer contacto (C3) & 18:30:35.9 & 20:30:35 & 4.64 & 285.8 & 1.0000 & 1.0000 \\
Cuarto contacto (C4) & 19:21:19.2 & 21:21:19 & -4.48 & 294.1 & 0.0000 & 0.0000 \\
\bottomrule
\end{tabular}
\end{table}
